\itemhead{1.2 例题 · 例题4（教材p.17）}{XJ-B1-C01-EX-P017-Y1338}
\label{XJ-B1-C01-EX-P017-Y1338}
{\footnotesize 湘教版教材p.17 / PDF第24页 / 1.2 例题 / 例题4}
\field{原题}{从“充分而不必要条件”“必要而不充分条件”“充要条件”和“既不充分又不必要条件”中选择适当的一种填空。
\subq{1}$a\ge5$ 是 $a$ 为正数的\blank；
\subq{2}四边形的两对角线相等是该四边形为矩形的\blank；
\subq{3}四边形的一组对边平行且相等是四边形的两组对边分别平行的\blank；
\subq{4}若 $x\in\mathbb R$，则 $x^2=2$ 是 $x=2$ 的\blank。}
\field{参考答案}{依次为充分而不必要、必要而不充分、充要、既不充分又不必要条件。}
{\small \field{关联知识}{L02 充分、必要、充要条件（教材帮p.30、31 / PDF39、40）}
\field{方法与易错点}{分别检查两个方向；用解集包含或具体反例判定，不凭句子位置猜方向。 易错：把“必要”箭头反写；只有一个方向证据便断言“不必要”；忽略给定论域。}
\field{答案依据}{例题原解位置：湘教版教材p.17及题库原图；独立编写可编辑参考解答。答案册此部分主要收练习、习题，不伪造例题答案册位置。}}
